\documentclass[12pt]{article}
\usepackage[utf8]{inputenc}
\usepackage[T1]{fontenc}
\usepackage{hyperref}
\title{Poqomam, Eastern Language Resources for K-12 Curriculum Development}
\author{Compiled by Claude Sonnet 5.5 (claude-sonnet-5-5) for Andr\'as Kornai}
\date{}
\begin{document}
\maketitle

\section{Introduction}
Poqomam, Eastern (ISO 639-3 poa; Glottocode poqo1257) is a retired code: it was merged into Poqomam [poc] effective 2009-01-16, because it denoted a dialect group. The prompt's ``family: Bookkeeping'' reflects Glottolog's label for retired entries, not a real language family; Poqomam is a Mayan language of Guatemala, with its largest pocket in Jalapa department. Ethnologue's 2009 edition gave 12,500 speakers for poa (1990). Resources below therefore concern Poqomam in general.

\section{Primary Educational Resources}
\begin{itemize}
\item News report on the ALMG virtual library for learning Poqomam (bilingual reading and writing guides, teacher materials, the Kumatzij text for second-language teaching): \url{https://agn.gt/sabias-que-existe-una-biblioteca-virtual-para-aprender-poqomam/}
\item News report on Mineduc training 829 teachers in Poqomam: \url{https://agn.gt/mineduc-concluye-capacitacion-de-829-docentes-en-idioma-poqomam/}
\end{itemize}

\section{Linguistic and Reference Materials}
\begin{itemize}
\item Glottolog retired entry for Poqomam, Eastern: \url{https://glottolog.org/resource/languoid/id/poqo1257}
\item Glottolog entry for Poqomam: \url{https://glottolog.org/resource/languoid/id/poqo1253}
\item Poqomam language (Wikipedia): \url{https://en.wikipedia.org/wiki/Poqomam_language}
\item News report on the ALMG bilingual Poqomam-Spanish dictionary (San Luis Jilotepeque, Jalapa, among the base communities): \url{https://agn.gt/?p=402251}
\end{itemize}

\section{Religious and Community Resources}
The Academia de Lenguas Mayas de Guatemala (ALMG) Poqomam linguistic community is the relevant body; no direct URL for it appeared in the search results.

\section{Digital and Audio Resources}
The news report above describes the multidialectal dictionary DALIMA and the Poqomam sound atlas Alispoq, useful for comparing dialects including Jalapa. Direct links were not seen in search results.
\begin{itemize}
\item Endangered Languages Project entry: \url{https://endangeredlanguages.com/lang/4986}
\end{itemize}

\section{Curriculum Development Context}
Poqomam has an official ALMG writing standard, a bilingual dictionary and teacher training in Guatemala's bilingual education system. A curriculum for the Jalapa variety should use the general Poqomam materials and consult the ALMG about dialect differences.
\end{document}